\scanpage{141}
\begin{center}- 129 -\end{center}
\margin{n\textsuperscript{o} 316}

où, dans la série $o_3$, tous les monômes sont de degré $\geq3$, par rapport à l'ensemble des $Y_i$ et des $Z_i$. En déduire que :
\[
\begin{aligned}
(3)\quad u(f(x,y))-u(f(y,x))&=\sum_{i<j}(X_iY_j-X_jY_i)([T_i,T_j]^{(e)}(u))+\\
&\quad+o'_3(x,y)
\end{aligned}
\]
où tous les termes de $o'_3$ sont de degré total $\geq3$. Montrer que si $G$ est commutatif, $\unsure{g|\mathfrak g}$ est commutative.

d) Soit $G'$ un second groupe formel, de dimension $n$, dont la loi de groupe est $f'$. On appelle \uline{homomorphisme formel} de $G$ dans $G'$ un système $F=(F_j(x))_{1\leq j\leq m}$ d'éléments de $A_x$ tel que l'on ait $f'(F(x),F(y))=F(f(x,y))$. Pour toute dérivation invariante à gauche $D\in\unsure{g|\mathfrak g}$, montrer que $v\longrightarrow D(v\circ F)$ est une dérivation invariante à gauche appartenant à l'algèbre de Lie $\unsure{g'|\mathfrak g'}$ de $G'$. Si on la désigne par $F^*(D)$, $F^*$ est un homomorphisme (un p-homomorphisme pour $p>0$) de $\unsure{g|\mathfrak g}$ dans $\unsure{g'|\mathfrak g'}$. Si $(T_i)_{1\leq i\leq n}$, $(T'_j)_{1\leq j\leq m}$ sont les bases de $\unsure{g|\mathfrak g}$ et $\unsure{g'|\mathfrak g'}$ telles que $T_i^{(e)}=D_i^{(e)}$ et ${T'_j}^{(e)}={D'_j}^{(e)}$, montrer que la matrice de $F^*$, rapportée à ces bases, est la matrice $(\partial F_j/\partial X_i)_c$ (terme constant de $\partial F_j/\partial X_i$, où $i$ est l'indice des lignes, $j$ celui des colonnes).

25) - On appelle \uline{groupe linéaire formel} à $n$ variables sur le corps $K$, et on désigne par $\mathrm{GL}(n,K)$ (lorsqu'aucune confusion n'est possible) le groupe formel de dimension $n^2$ dont la loi de groupe est définie par
\[
f_{ij}(u,v)=-\delta_{ij}+\sum_{k=1}^n(\delta_{ik}+U_{ik})(\delta_{kj}+V_{kj})
\]
($\delta_{ij}$ indice de Kronecker ; les $3n^2$ indéterminées de la définition générale de l'exerc. 24 sont ici notées $U_{ij}$, $V_{ij}$, $W_{ij}$ avec $2$ indices variant de $1$ à $n$). Montrer que, si $D_{ij}=\partial/\partial U_{ij}$, les dérivations
\margin{140}
